% Options for packages loaded elsewhere
\PassOptionsToPackage{unicode}{hyperref}
\PassOptionsToPackage{hyphens}{url}
\documentclass[
  japanese,
]{ltjsarticle}
\usepackage{xcolor}
\usepackage[margin=25mm]{geometry}
\usepackage{amsmath,amssymb}
\setcounter{secnumdepth}{-\maxdimen} % remove section numbering
\usepackage{iftex}
\ifPDFTeX
  \usepackage[T1]{fontenc}
  \usepackage[utf8]{inputenc}
  \usepackage{textcomp} % provide euro and other symbols
\else % if luatex or xetex
  \usepackage{unicode-math} % this also loads fontspec
  \defaultfontfeatures{Scale=MatchLowercase}
  \defaultfontfeatures[\rmfamily]{Ligatures=TeX,Scale=1}
\fi
\usepackage{lmodern}
\ifPDFTeX\else
  % xetex/luatex font selection
\fi
% Use upquote if available, for straight quotes in verbatim environments
\IfFileExists{upquote.sty}{\usepackage{upquote}}{}
\IfFileExists{microtype.sty}{% use microtype if available
  \usepackage[]{microtype}
  \UseMicrotypeSet[protrusion]{basicmath} % disable protrusion for tt fonts
}{}
\makeatletter
\@ifundefined{KOMAClassName}{% if non-KOMA class
  \IfFileExists{parskip.sty}{%
    \usepackage{parskip}
  }{% else
    \setlength{\parindent}{0pt}
    \setlength{\parskip}{6pt plus 2pt minus 1pt}}
}{% if KOMA class
  \KOMAoptions{parskip=half}}
\makeatother
% definitions for citeproc citations
\NewDocumentCommand\citeproctext{}{}
\NewDocumentCommand\citeproc{mm}{%
  \begingroup\def\citeproctext{#2}\cite{#1}\endgroup}
\makeatletter
 % allow citations to break across lines
 \let\@cite@ofmt\@firstofone
 % avoid brackets around text for \cite:
 \def\@biblabel#1{}
 \def\@cite#1#2{{#1\if@tempswa , #2\fi}}
\makeatother
\newlength{\cslhangindent}
\setlength{\cslhangindent}{1.5em}
\newlength{\csllabelwidth}
\setlength{\csllabelwidth}{3em}
\newenvironment{CSLReferences}[2] % #1 hanging-indent, #2 entry-spacing
 {\begin{list}{}{%
  \setlength{\itemindent}{0pt}
  \setlength{\leftmargin}{0pt}
  \setlength{\parsep}{0pt}
  % turn on hanging indent if param 1 is 1
  \ifodd #1
   \setlength{\leftmargin}{\cslhangindent}
   \setlength{\itemindent}{-1\cslhangindent}
  \fi
  % set entry spacing
  \setlength{\itemsep}{#2\baselineskip}}}
 {\end{list}}
\usepackage{calc}
\newcommand{\CSLBlock}[1]{\hfill\break\parbox[t]{\linewidth}{\strut\ignorespaces#1\strut}}
\newcommand{\CSLLeftMargin}[1]{\parbox[t]{\csllabelwidth}{\strut#1\strut}}
\newcommand{\CSLRightInline}[1]{\parbox[t]{\linewidth - \csllabelwidth}{\strut#1\strut}}
\newcommand{\CSLIndent}[1]{\hspace{\cslhangindent}#1}
\ifLuaTeX
\usepackage[bidi=basic,shorthands=off,,provide=*]{babel}
\else
\usepackage[bidi=default,shorthands=off,,provide=*]{babel}
\fi
\ifLuaTeX
  \usepackage{selnolig} % disable illegal ligatures
\fi
\setlength{\emergencystretch}{3em} % prevent overfull lines
\providecommand{\tightlist}{%
  \setlength{\itemsep}{0pt}\setlength{\parskip}{0pt}}
\usepackage{bookmark}
\IfFileExists{xurl.sty}{\usepackage{xurl}}{} % add URL line breaks if available
\urlstyle{same}
\hypersetup{
  pdftitle={レート制御とモニタの形式化},
  pdflang={ja-JP},
  hidelinks,
  pdfcreator={LaTeX via pandoc}}

\title{レート制御とモニタの形式化}
\usepackage{etoolbox}
\makeatletter
\providecommand{\subtitle}[1]{% add subtitle to \maketitle
  \apptocmd{\@title}{\par {\large #1 \par}}{}{}
}
\makeatother
\subtitle{DeepLbot3を事例とした研究ノート}
\author{}
\date{2026年10月28日}

\begin{document}
\maketitle

\section{目的と対象}\label{ux76eeux7684ux3068ux5bfeux8c61}

DeepLbot3のDeepL呼び出しを事例に、送信制御と実行時観測を研究する。ここで送信はアプリケーションによる
\texttt{fetch}
呼び出しを表し、DeepLへの到着時刻ではない。バケットのパラメータは研究者が定める方針であり、DeepL内部の制限の推定ではない。流量制約はネットワーク計算
(Boudec と Thiran 2001; Boudec 2018)、実行時の観測は実行時検証 (Bartocci
ほか 2018) に関連する。

以下は有限実行を対象とする数学的説明である。各理論はIsabelle-checkerの
\texttt{strict=true}
で検査済みだが、TypeScript実装との対応は未証明である。

\section{トークンバケット}\label{ux30c8ux30fcux30afux30f3ux30d0ux30b1ux30c3ux30c8}

容量 \(B\)、補充率 \(r\)、残量 \(b\)、時間増分 \(d\)
は自然数とする。\(d=0\) も許す。イベントは時間経過
\(\mathrm{Tick}(d)\)、送信 \(\mathrm{Send}\)、その他 \(\mathrm{Other}\)
とし、状態遷移を以下で定める。

\[
\begin{aligned}
b &\xrightarrow{\mathrm{Tick}(d)} \min(B,b+rd),\\
b &\xrightarrow{\mathrm{Send}} b-1 && (b>0),\\
b &\xrightarrow{\mathrm{Other}} b.
\end{aligned}
\]

有限列 \(E\) の送信回数を \(S(E)\)、時間経過の和を \(T(E)\)
とする。\(b\xrightarrow{E}b'\)
は列の各イベントが順に遷移可能であることを意味する。これは
\texttt{Rate\_Bucket.thy} の
\texttt{step}、\texttt{run}、\texttt{sends}、\texttt{elapsed}
に対応する。

\textbf{補題1（残量上限）}　\(b\le B\) かつ \(b\xrightarrow{E}b'\)
ならば \(b'\le B\)。

\emph{証明.} 実行列について帰納法を用いる。空列では明らか。時間経過では
\(\min\) が \(B\)
を超えず、送信では残量が減り、その他のイベントでは変わらない。よって各遷移で上限が保存される。□

\textbf{補題2（資源収支）}　\(b\xrightarrow{E}b'\) ならば
\(S(E)+b'\le b+rT(E)\)。

\emph{証明.}
列の先頭からの帰納法。空列では等号。送信では送信回数が1増え残量が1減るので左辺は変わらない。時間経過
\(d\) による残量の増加は高々 \(rd\)
である。その他のイベントはどちらも変えない。したがって不等式は保存される。□

\textbf{定理1（送信量上限）}　\(b\le B\) かつ \(b\xrightarrow{E}b'\)
ならば \(S(E)\le B+rT(E)\)。

\emph{証明.} 補題2と \(b'\ge0\) より \(S(E)\le b+rT(E)\)。仮定
\(b\le B\) を適用する。□

任意の有限実行区間に適用できる。複数実行のイベントが一つの共有バケットの列に直列化できれば、この定理を使えるが、直列化と共有状態の実装保証はまだ含まれない。

\section{実行可能モニタ}\label{ux5b9fux884cux53efux80fdux30e2ux30cbux30bf}

モニタ状態は \(\mathrm{Some}(b)\) または \(\mathrm{None}\)
とする。前者ではバケットと同じ補充・消費を行い、残量0で送信を観測すると後者へ移る。後者は以後維持される。列
\(E\) の処理結果を \(M(B,r,b,E)\) とする。実装は
\texttt{Rate\_Monitor.thy} にある。

\textbf{定理2（遷移関係との等価性）}　\(b\xrightarrow{E}b'\) と
\(M(B,r,b,E)=\mathrm{Some}(b')\) は同値。

\emph{証明.}
単一イベントについて、遷移が存在することとモニタ更新が同じ残量を返すことは、3種類のイベントの定義から同値である。列の長さで帰納する。空列は自明。先頭イベントに単一イベントの同値性、残りの列に帰納法の仮定を用いる。途中で
\(\mathrm{None}\) になると復帰しないので、成功結果にはならない。□

\textbf{定理3（受理列の送信量）}　\(b\le B\) かつ
\(M(B,r,b,E)=\mathrm{Some}(b')\) ならば \(S(E)\le B+rT(E)\)。

\emph{証明.} 定理2でバケットの実行関係を得て、定理1を適用する。□

ここで証明した等価性は、同じ操作的バケット意味論との対応である。独立した到着曲線仕様との同値性や、イベントの欠落がない観測はまだ証明していない。

\section{ボットの処理段階}\label{ux30dcux30c3ux30c8ux306eux51e6ux7406ux6bb5ux968e}

単一実行の段階を準備、送信可能、要求中、応答検査、返信、利用記録、完了リアクション、終了とする。成功経路はこの順序で進み、失敗遷移は任意の段階から終了する。返信成功を
\(\mathrm{ReplyOK}\)、完了リアクション成功を \(\mathrm{MarkOK}\)
と表す。これは実装のすべての分岐を表すものではない。理論は
\texttt{Bot\_Order.thy} にある。

\textbf{補題3}　任意の段階 \(p\)
からの実行が完了リアクション段階へ到達するか、実行列に
\(\mathrm{MarkOK}\) を含むなら、\(p\)
は利用記録または完了リアクション段階であるか、実行列に
\(\mathrm{ReplyOK}\) が含まれる。

\emph{証明.}
実行列について帰納する。空列の場合、結論は開始段階の条件から従う。先頭遷移を場合分けすると、利用記録に入る唯一の成功遷移は返信成功、完了リアクションに入る唯一の遷移は利用記録終了であり、他の遷移は帰納法の仮定を保存する。失敗遷移は終了段階に進む。よって各場合で成立する。□

\textbf{定理4（返信成功の存在）}　準備状態からの有限実行列に
\(\mathrm{MarkOK}\) が含まれるならば、同じ列に \(\mathrm{ReplyOK}\)
が含まれる。

\emph{証明.}
補題3を適用する。開始段階は準備なので、利用記録・完了リアクションであるという例外は成立しない。□

現在のIsabelle定理はイベントの存在を示す。列中の厳密な先行順序を直接述べる定理は未実装である。

\section{限界と今後の作業}\label{ux9650ux754cux3068ux4ecaux5f8cux306eux4f5cux696d}

実行IDと共有バケットの並行合成、再試行・期限、計装の消去、観測の完全性、許可取得と呼び出しの時間差、TypeScriptとの精密化関係は未形式化である。線形化可能性
(Herlihy と Wing 1990) と時間オートマトンの検証 (Wimmer と Lammich 2018)
は今後の検討に関連する。Slack環境の実測は仮定の妥当性を評価するが、観測された最大遅延は全実行の証明済み上限ではない。

\section*{参考文献}\label{ux53c2ux8003ux6587ux732e}
\addcontentsline{toc}{section}{参考文献}

\protect\phantomsection\label{refs}
\begin{CSLReferences}{1}{0}
\bibitem[\citeproctext]{ref-runtime}
Bartocci, Ezio, Ylies Falcone, Adrian Francalanza, と Giles Reger. 2018.
{「Introduction to Runtime Verification」}. \emph{Lectures on Runtime
Verification}. Springer.
\url{https://doi.org/10.1007/978-3-319-75632-5_1}.

\bibitem[\citeproctext]{ref-regulators}
Boudec, Jean-Yves Le. 2018. {「A Theory of Traffic Regulators for
Deterministic Networks With Application to Interleaved Regulators」}.
\emph{IEEE/ACM Transactions on Networking} 26 (6): 2721--33.
\url{https://doi.org/10.1109/TNET.2018.2875191}.

\bibitem[\citeproctext]{ref-network}
Boudec, Jean-Yves Le, と Patrick Thiran. 2001. \emph{Network Calculus: A
Theory of Deterministic Queuing Systems for the Internet}. Springer.
\url{https://doi.org/10.1007/3-540-45318-0}.

\bibitem[\citeproctext]{ref-linear}
Herlihy, Maurice, と Jeannette M. Wing. 1990. {「Linearizability: A
Correctness Condition for Concurrent Objects」}. \emph{ACM Transactions
on Programming Languages and Systems} 12 (3): 463--92.
\url{https://doi.org/10.1145/78969.78972}.

\bibitem[\citeproctext]{ref-timed}
Wimmer, Simon, と Peter Lammich. 2018. {「Verified Model Checking of
Timed Automata」}. \emph{TACAS}, 61--78. Springer.
\url{https://doi.org/10.1007/978-3-319-89960-2_4}.

\end{CSLReferences}

\end{document}
